% Problema 7, Capítulo 2 (Spivak)
% Sumas de potencias como polinomios

\section*{Problema 7 - Fórmula general para $\sum_{k=1}^n k^p$}
\addcontentsline{toc}{section}{Problema 7 - Fórmula general para $\sum_{k=1}^n k^p$}

El objetivo es mostrar que para cada entero positivo $p$ existe un polinomio
en $n$ de grado $p+1$ con coeficiente dominante $1/(p+1)$ tal que
\[
\sum_{k=1}^n k^p = \frac{n^{p+1}}{p+1} + A n^p + B n^{p-1}+\cdots
t\]
donde los demás coeficientes son racionales que dependen sólo de $p$.

La demostración se basa en el método del problema anterior: escribir
a partir de la identidad
\[
(k+1)^{p+1}-k^{p+1} = (p+1)k^p + \binom{p+1}{2}k^{p-1} + \cdots + 1
t\]
y sumar telescópicamente de $k=1$ hasta $n$. Al desplazar términos obtenemos
a una ecuación lineal entre las sumas de potencias de órdenes $p,p-1,\dots,0$.

Esquema general:
\begin{align*}
(n+1)^{p+1}-1 &= \sum_{k=1}^n\big((k+1)^{p+1}-k^{p+1}\big) \\
&= (p+1)\sum_{k=1}^n k^p + \binom{p+1}{2}\sum_{k=1}^n k^{p-1} + \cdots + \sum_{k=1}^n1.
\end{align*}

Denotemos $S_j(n)=\sum_{k=1}^n k^j$. La igualdad anterior se reescribe como
\[
(p+1)S_p(n) = (n+1)^{p+1}-1 - \binom{p+1}{2}S_{p-1}(n) - \cdots - S_0(n).
\]
Por inducción en $p$, suponemos que cada $S_j(n)$ con $j<p$ es un polinomio de
grado $j+1$ cuyos coeficientes conocidos se han calculado anteriormente.
Entonces el lado derecho es un polinomio en $n$ de grado $p+1$ cuyo
coeficiente líder es precisamente $1$ (viene de $(n+1)^{p+1}$). Dividiendo
por $p+1$ obtenemos el polinomio buscado para $S_p(n)$ con coeficiente
líder $1/(p+1)$.

\subsection*{Ejemplos concretos}

Las primeras diez fórmulas obtenidas por esta técnica son (tal como se
muestra en la imagen del enunciado original):
\begin{align*}
\sum_{k=1}^n k &= \tfrac12 n^2 + \tfrac12 n, \\
\sum_{k=1}^n k^2 &= \tfrac13 n^3 + \tfrac12 n^2 + \tfrac16 n, \\
\sum_{k=1}^n k^3 &= \tfrac14 n^4 + \tfrac12 n^3 + \tfrac14 n^2 = \left(\tfrac12 n^2 + \tfrac12 n\right)^2, \\
\sum_{k=1}^n k^4 &= \tfrac15 n^5 + \tfrac12 n^4 + \tfrac13 n^3 - \tfrac{1}{30}n, \\
\sum_{k=1}^n k^5 &= \tfrac16 n^6 + \tfrac12 n^5 + \tfrac56 n^4 - \tfrac16 n^2, \\
\sum_{k=1}^n k^6 &= \tfrac17 n^7 + \tfrac12 n^6 + \tfrac56 n^5 - \tfrac17 n^3 + \tfrac16 n, \\
\sum_{k=1}^n k^7 &= \tfrac18 n^8 + \tfrac12 n^7 + \tfrac78 n^6 - \tfrac17 n^4 + \tfrac17 n^2, \\
\sum_{k=1}^n k^8 &= \tfrac19 n^9 + \tfrac12 n^8 + \tfrac89 n^7 - \tfrac78 n^5 + \tfrac89 n^3 - \tfrac{1}{30}n, \\
\sum_{k=1}^n k^9 &= \tfrac1{10} n^{10} + \tfrac12 n^9 + \tfrac89 n^8 - \tfrac78 n^6 + \tfrac19 n^4 - \tfrac3{10} n^2, \\
\sum_{k=1}^n k^{10} &= \tfrac1{11} n^{11} + \tfrac12 n^{10} + \tfrac9{10} n^9 - n^7 + \tfrac12 n^5 - \tfrac1{2}n^3 + \tfrac5{66}n.
\end{align*}

Observar que el segundo coeficiente es siempre $\tfrac12$ y después de la
tercera columna las potencias aparecen con coeficientes decrecientes de dos
en dos. La demostración precedente muestra que estos patrones se derivan de
las binomiales en la identidad telescópica.

\subsection*{Consecuencias y generalizaciones}

En particular, la existencia de dicho polinomio implica que
$S_p(n)=O(n^{p+1})$ y que
\[
\frac{S_p(n)}{n^{p+1}}\to \frac{1}{p+1} \quad(n\to\infty),
\]
una forma elemental de establecer el comportamiento asintótico de las
sumas de potencias. Este argumento es la base de la llamada "fórmula de
Faulhaber" y conecta con los números de Bernoulli cuando se busca una
expresión explícita para los coeficientes.

\vspace{1em}
Puedes probar inductivamente el paso anterior y calcular coeficientes
determinados mediante comparaciones de términos en ambos lados de la identidad
telecópica.

